\begin{abstract}
Unsymmetric (Petrov-Galerkin) spectral elements use conforming parametric test functions and metric trial functions that are complete in the physical coordinates. On meshes whose edge curvature persists under refinement, they therefore converge with the optimal rate, whereas the conventional spectral element of degree $p$ with a quadratic geometry is complete only to degree $\lfloor p/2\rfloor$. The metric trial functions are, however, obtained by nodal interpolation in the tensor-product polynomials of skew coordinates, and on strongly distorted elements they overshoot between the nodes, increasingly with the degree. In the paper at hand, we separate two causes of this behavior and propose one remedy for each. The first cause is the conditioning of the construction. Since the metric shape functions do not depend on the polynomial basis of the generalized Vandermonde matrix, we use the Lagrange polynomials through the reference Gauss-Lobatto-Legendre (GLL) points in skew coordinates scaled to the bounding box of the nodes. On parallelograms, the Vandermonde matrix is then a permutation matrix, and on distorted elements its condition number at $p=16$ is three to four orders of magnitude smaller than with monomials. Its condition number equals the Lebesgue function of the metric interpolation at the reference grid of the bounding box, so that no basis can reduce the Lebesgue constant itself. The second cause is the interpolation problem. Since exact reproduction requires only the complete polynomials of degree $p$ in the physical coordinates, we write only these in skew coordinates and the remaining functions, the $L^2$-orthogonal complement of the complete polynomials in the tensor-product space, in the reference coordinates, in which the nodes form the exact GLL grid. The resulting hybrid element coincides with the conventional spectral element on every straight-sided quadrilateral, and its unisolvence reduces to the regularity of a matrix of the size of the complete polynomial space. Its Lebesgue constant stays at or close to the GLL value for moderate edge curvature, where that of the tensor-product metric space reaches $10^4$--$10^7$ at $p=16$. On meshes with persistent edge curvature, the hybrid element converges with the optimal rate for $p\le6$, whereas the original unsymmetric element diverges for $p=5,6$. For a distorted thick cylinder, its error is up to 390 times smaller, and it is 0.78 to 1.00 times the interpolation error of its trial space on all meshes without folded elements. Moreover, the eigenvalues with negative real part of the original element are removed, although a slow growth of the transient response remains. The price is an error that is 1.3 to 5.5 times larger where the original element is stable and, on folded elements, singular integrands that limit the attainable accuracy.
\end{abstract}
